%% ---------------------------------------------------------------------
\section{Estabilidad desde la malla abierta: márgenes de ganancia y de fase}
\label{sec:rf-estabilidad}
%% ---------------------------------------------------------------------

\subsection{La frontera: $|\Gma|=1$ y fase $-180\grados$}

La ecuación característica del lazo cerrado es $1+\Gma(s)=0$, es decir, $\Gma(s)=-1$. Evaluada sobre el eje imaginario, la frontera entre estabilidad e inestabilidad es una frecuencia en la que, a la vez,
\begin{equation}
  |\Gma(\jw)|=1\ \ (0\dB)\qquad\text{y}\qquad \angle\Gma(\jw)=-180\grados .
  \label{eq:rf-frontera}
\end{equation}
La intuición es la de la \cref{sec:rf-lenguaje}: una senoide que da la vuelta al lazo regresa del mismo tamaño ($|\Gma|=1$) e invertida ($-180\grados$); el signo menos del comparador la vuelve a invertir, se suma a sí misma y el lazo oscila solo. Si en la frecuencia donde la fase vale $-180\grados$ la senoide regresa \emph{más pequeña} ($|\Gma|<1$), cada vuelta la achica y la oscilación muere.

\begin{ideaclave}
Si la malla abierta es estable y de fase mínima, el lazo cerrado es asintóticamente estable cuando la magnitud de $\Gma(\jw)$ es menor que $0\dB$ en la frecuencia en que su fase vale $-180\grados$.
\end{ideaclave}

\subsection{Márgenes de ganancia y de fase}

Saber que el lazo es estable no basta: queremos saber \emph{qué tan lejos} está de la frontera. Los dos márgenes miden esa distancia, cada uno en una dirección (\cref{fig:rf-margenes}).

\begin{definicion}[label=def:rf-gm]{Margen de ganancia}
El \clave{margen de ganancia} ($\GM$) es cuánto se puede aumentar la ganancia de la malla abierta, en decibeles, en la frecuencia donde la fase vale $-180\grados$ ($\wgm$, cruce de fase), antes de que el lazo cerrado se vuelva inestable:
\begin{equation}
  \GM=-20\log|\Gma(\jw_{GM})|\quad\text{dB},\qquad \angle\Gma(\jw_{GM})=-180\grados.
\end{equation}
\end{definicion}

\begin{definicion}[label=def:rf-pm]{Margen de fase}
El \clave{margen de fase} ($\PM$) es cuánto atraso de fase adicional se tolera en la frecuencia donde $|\Gma|=1$ ($\wpm$, cruce de ganancia) antes de llegar a $-180\grados$:
\begin{equation}
  \PM=180\grados+\angle\Gma(\jw_{PM}),\qquad |\Gma(\jw_{PM})|=1\ (0\dB).
\end{equation}
\end{definicion}

\begin{figure}[ht]
\centering
\includegraphics[width=0.88\linewidth]{cap10/f_margenes_def}
\caption{Lectura de los márgenes. El $\GM$ se mide en magnitud, en la frecuencia donde la fase cruza $-180\grados$; el $\PM$ se mide en fase, en la frecuencia donde la magnitud cruza $0\dB$.}
\label{fig:rf-margenes}
\end{figure}

\begin{cuidado}[Los nombres de MATLAB invitan a confundirse]
En \cmd{[Gm,Pm,Wcg,Wcp] = margin(G)}, \cmd{Wcg} (\emph{gain crossover} en otros textos) es en realidad la frecuencia donde se mide el $\GM$, es decir, nuestro $\wgm$, y \cmd{Wcp} es la frecuencia del $\PM$, nuestro $\wpm$. Además \cmd{Gm} sale en \emph{veces}: el $\GM$ en dB es \cmd{20*log10(Gm)}. La función del curso \cmd{margenesTabla(G)} imprime todo ya convertido.
\end{cuidado}

\begin{ejemplo}[label=ej:rf-margenes]{Márgenes asintóticos contra exactos}
Determine el $\GM$ y el $\PM$ de $\Gma(s)=\dfrac{10\,000}{(s+1)(s+10)(s+100)}$ con el diagrama asintótico y compare con \cmd{margin}.
\solucion \textbf{Forma de Bode:} $\Gma(s)=\dfrac{10}{(s+1)(s/10+1)(s/100+1)}$, es decir, $20\dB$ en bajas frecuencias y tres esquinas.

\textbf{Cruce por $0\dB$ (PM).} La asíntota baja $-20\dBdec$ desde $\omega=1$ y llega a $0\dB$ justo en $\omega=10$. La fase asintótica allí: el polo en $1$ ya aportó $-90\grados$, el de $10$ está en su esquina ($-45\grados$) y el de $100$ todavía no empieza. Total $-135\grados$: $\PM\approx45\grados$.

\textbf{Cruce por $-180\grados$ (GM).} Entre $\omega=10$ y $\omega=100$ la fase asintótica cae $-45\grados$ por década por el polo en $10$ y otros $-45\grados$ por década por el de $100$, desde $-135\grados$: llega a $-180\grados$ media década después de $10$, en $\wgm=10^{1.5}=31.6$\,rad/s. La magnitud asintótica allí es $0-40\times0.5=-20\dB$: $\GM\approx20\dB$.

\textbf{Comparación.}
\begin{center}\small
\begin{tabular}{@{}lcc@{}}
\toprule
 & asintótico & exacto (\cmd{margin})\\
\midrule
$\PM$ & $45\grados$ ($\omega=10$) & $54.9\grados$ ($\omega=7.80$)\\
$\GM$ & $20\dB$ ($\omega=31.6$) & $21.7\dB$ ($\omega=33.3$)\\
\bottomrule
\end{tabular}
\end{center}
El asintótico es conservador aquí: sirve para estimar en el tablero; el número final siempre con \cmd{margin}.
\end{ejemplo}

\begin{figure}[ht]
\centering
\includegraphics[width=0.9\linewidth]{cap10/f_margenes_asint}
\caption{Ejemplo~\ref{ej:rf-margenes}: los cuatro factores, su suma asintótica y la curva real, con los márgenes asintóticos marcados.}
\label{fig:rf-margenes-asint}
\end{figure}

\codigorepo{cap10_respuesta_frecuencial/Clase2_Margenes/C2_Ej2_Margenes.m}{margenes}{ejemplo \ref{ej:rf-margenes}}
\video{V3_asintotico_margenes}{El Bode se arma término a término; al final, PM y GM sobre la asintótica.}

%% ---------------------------------------------------------------------
\subsection{¿Qué significan los márgenes en la práctica?}
\label{sec:rf-significado}

\paragraph{El GM: el micrófono y el parlante}
\notamargen{El ``aullido'' de un sistema de sonido es un lazo cerrado que perdió su margen de ganancia.}
Un micrófono frente a un parlante forma un lazo: lo que sale del parlante vuelve a entrar al micrófono. Subir el volumen es subir la ganancia de la malla abierta \emph{sin cambiar la fase}. Mientras haya margen de ganancia no pasa nada; cuando el volumen sube más que el $\GM$, en la frecuencia donde la fase acústica vale $-180\grados$ aparece el aullido: una oscilación sostenida. El $\GM$ es, literalmente, cuánto se puede subir el volumen.

\paragraph{El PM: un presupuesto de fase para lo que no modelamos}
Ningún modelo es perfecto. Sensores, actuadores y filtros que dejamos por fuera agregan atraso de fase que la malla abierta modelada no ve. El $\PM$ es el \emph{presupuesto} con el que pagamos ese atraso. Para verlo, pongamos un sensor de primer orden $\Hs(s)=a/(s+a)$ en la retroalimentación del ejemplo~\ref{ej:rf-margenes} (\cref{fig:rf-lazosensor}). En el cruce, el sensor resta aproximadamente $\tan^{-1}(\wpm/a)$, de modo que
\begin{equation}
  \PM_{\text{nuevo}}\approx\PM-\tan^{-1}\frac{\wpm}{a}\qquad(\text{buena estimación si } a\gg\wpm).
\end{equation}

\begin{figure}[ht]
\centering
\includegraphics[width=0.55\linewidth]{cap10/lazo_sensor}
\caption{El lazo del ejemplo~\ref{ej:rf-margenes} con un sensor de primer orden. La malla abierta es ahora $\Gma=G\,\Hs$.}
\label{fig:rf-lazosensor}
\end{figure}

\begin{center}\small
\begin{tabular}{@{}ccccc@{}}
\toprule
$a$ (rad/s) & $\PM$ estimado & $\PM$ (\cmd{margin}) & $\GM$ & sobrepico\\
\midrule
sin sensor & -- & $54.9\grados$ & $21.7\dB$ & $18\%$\\
$50$ & $46.0\grados$ & $46.5\grados$ & $12.7\dB$ & $28\%$\\
$20$ & $33.6\grados$ & $36.4\grados$ & $8.6\dB$ & $44\%$\\
$10$ & $17.0\grados$ & $26.3\grados$ & $6.1\dB$ & $68\%$\\
$5$  & $-2.4\grados$ & $18.7\grados$ & $4.7\dB$ & $104\%$\\
\bottomrule
\end{tabular}
\captionof{table}{El sensor ``gasta'' fase. Con $a=50$, casi diez veces más rápido que el cruce, la estimación es buena; con sensores lentos la estimación es pesimista (el sensor también baja la magnitud y corre el cruce a la izquierda), pero la tendencia es la misma.}
\label{tab:rf-sensor}
\end{center}

\begin{figure}[ht]
\centering
\includegraphics[width=\linewidth]{cap10/f_sensor}
\caption{Un sensor más lento gasta más fase en el cruce (derecha) y el lazo cerrado se vuelve más oscilatorio (izquierda), aunque la planta sea la misma.}
\label{fig:rf-sensor}
\end{figure}

\codigorepo{cap10_respuesta_frecuencial/Clase2_Margenes/C2_Video_SensorMargenFase.m}{sensor}{tabla \ref{tab:rf-sensor}}
\video{V12_sensor_margen_fase}{El polo del sensor baja de $300$ a $5$\,rad/s: se ve cómo se consume el PM y crece el sobrepico.}

\begin{ideaclave}
El $\GM$ dice cuánta ganancia de más aguanta el lazo; el $\PM$, cuánto atraso de fase no modelado aguanta. Un buen diseño necesita los dos.
\end{ideaclave}

\begin{profundiza}[Los márgenes como distancia al punto $-1$]
Ambos márgenes miden qué tan cerca pasa $\Gma(\jw)$ del punto crítico $-1$, cada uno en una dirección: el $\GM$ sobre el eje real negativo y el $\PM$ sobre el círculo unitario. Una medida que los resume es la distancia mínima $\min_\omega|1+\Gma(\jw)|=1/M_s$, donde $M_s=\max_\omega|\Sens(\jw)|$ es el pico de la sensitividad $\Sens=1/(1+\Gma)$. Un $M_s$ entre $1.3$ y $2$ garantiza a la vez $\GM\geq20\log\frac{M_s}{M_s-1}$ y $\PM\geq2\sin^{-1}\frac{1}{2M_s}$; por ejemplo, $M_s=2$ asegura $\GM\geq6\dB$ y $\PM\geq29\grados$~\citep{AsMu10}.
\end{profundiza}

%% ---------------------------------------------------------------------
\subsection{¿Cuánto margen es suficiente?}

\begin{itemize}
  \item \textbf{Fase mínima:} $\GM$ y $\PM$ positivos $\Rightarrow$ estable. Márgenes negativos $\Rightarrow$ inestable.
  \item \textbf{Regla práctica:} $30\grados<\PM<60\grados$ y $\GM>6\dB$ dan una estabilidad que tolera variaciones razonables de la ganancia y de las constantes de tiempo.
  \item \textbf{Fase no mínima} (ceros o polos en el semiplano derecho): el Bode solo no basta; se requiere un estudio más detallado (criterio de Nyquist).
\end{itemize}

\begin{pausa}
Tres lazos de fase mínima: (A) $\PM=40\grados$, $\GM=12\dB$; (B) $\PM=-15\grados$, $\GM=-4\dB$; (C) $\PM=70\grados$, $\GM=2\dB$. ¿Cuáles son estables? ¿Cuál le preocupa más?
\respuesta{(A) estable, con márgenes cómodos. (B) inestable: hay que bajar la ganancia al menos $4\dB$. (C) estable pero frágil: $2\dB$ es un factor $10^{2/20}=1.26$, así que un $26\%$ más de ganancia lo inestabiliza. Un $\PM$ grande no basta: hay que mirar los dos márgenes.}
\end{pausa}

%% ---------------------------------------------------------------------
\subsection{¿Cuánta ganancia aguanta? Routh-Hurwitz contra Bode}
\label{sec:rf-rangokc}

\begin{figure}[ht]
\centering
\includegraphics[width=0.62\linewidth]{cap10/lazo_basico}
\caption{Lazo de control con retroalimentación. En los ejemplos de diseño, $\Hs(s)=1$.}
\label{fig:rf-lazo}
\end{figure}

\begin{ejemplo}[label=ej:rf-rangokc]{Rango de $\Kc$ por dos caminos}
En el lazo de la \cref{fig:rf-lazo}, con $\Hs=1$, $\Gc(s)=\Kc$ y $\Gp(s)=\dfrac{1}{(s+1)(s+10)(s+100)}$, determine el rango de $\Kc$ que garantiza estabilidad asintótica con Routh-Hurwitz y con el diagrama de Bode.
\solucion \textbf{Routh-Hurwitz.} El polinomio característico es $s^3+111s^2+1110s+(1000+\Kc)$. La fila $s^1$ exige $111\cdot1110-(1000+\Kc)>0$ y la fila $s^0$, $1000+\Kc>0$:
\[
  -1000<\Kc<122\,210 .
\]
En $\Kc=\Kcr=122\,210$ el polinomio auxiliar $111s^2+123\,210=0$ da $s=\pm j\sqrt{1110}$: el lazo oscila con $\wcr=33.32$\,rad/s y periodo $\Pcr=2\pi/\wcr=0.189$\,s.

\textbf{Bode.} Para dibujar hay que fijar algún $\Kc$. Conviene el que deja la ganancia DC en $0\dB$, $\Kc=1000$, con lo que $\Gma(s)=\dfrac{1}{(s+1)(s/10+1)(s/100+1)}$. El cruce de fase asintótico está otra vez en $\omega=31.6$, donde la magnitud asintótica vale $-40\dB$: hay $\GM\approx+40\dB$, es decir, la ganancia se puede multiplicar por $100$. Así $\Kcr\approx1000\times100=100\,000$ y el rango positivo es $0<\Kc<100\,000$. Con \cmd{margin} el $\GM$ exacto es $41.74\dB$ ($\times122.2$), que reproduce $\Kcr=122\,210$.

La lectura asintótica se equivoca hacia el lado seguro: subestima la ganancia crítica.
\end{ejemplo}

\begin{figure}[ht]
\centering
\includegraphics[width=0.9\linewidth]{cap10/f_rango_kc}
\caption{Respuesta al escalón del lazo cerrado al subir $\Kc$: en $\Kc=\Kcr$ la oscilación ya no se amortigua.}
\label{fig:rf-rangokc}
\end{figure}

\codigorepo{cap10_respuesta_frecuencial/Clase2_Margenes/C2_Ej3_RangoKc.m}{routh}{Routh-Hurwitz sin Symbolic Toolbox}
\video{V4_barrido_Kc}{Subir $\Kc$ solo sube la magnitud: la barra roja del GM se acorta hasta cero en $\Kc=122\,210$.}

\begin{modelo}[LazoKc.slx]
\cmd{crearModelo\_LazoKc.m} arma el lazo escalón $\to$ suma $\to$ $\Kc$ $\to$ $\Gp(s)$ con retroalimentación unitaria. Pruebe $\Kc=1000$, $50\,000$, $122\,210$ (marginal) y $150\,000$ (inestable). Para otra planta cambie el parámetro \emph{Poles} del bloque $\Gp$.
\end{modelo}

\paragraph{Conexión con lo que ya vimos: Ziegler-Nichols}
\notamargen{Ziegler y Nichols publicaron sus reglas en 1942~\citep{ZiNi42}.}
El $\GM$ entrega la ganancia crítica y $\wgm$ la frecuencia crítica: son justo los datos del método de oscilación sostenida de Ziegler y Nichols (capítulo de sintonización). Con nuestro ejemplo, $\Kcr=122\,210$ y $\Pcr=0.189$\,s, y la tabla de ZN para un PID da
\[
  \Kc=0.6\,\Kcr=73\,326,\qquad \Ti=\Pcr/2=0.0943\ \text{s},\qquad \Td=\Pcr/8=0.0236\ \text{s}.
\]
La ventaja del Bode es que \emph{no hay que llevar la planta real al borde de la inestabilidad} para medir $\Kcr$ y $\Pcr$: se leen en el diagrama. En la \cref{sec:rf-pid} veremos qué margen de fase deja esa sintonización.
